\section*{अध्याय \ref{ch:g10:orders-of-magnitude} --- परिमाण की कोटि: ब्रह्मांड को मापना}

\begin{solution}{exo:g10:orders-of-magnitude:1}
\emph{1.} $\qty{750}{nm} = 750 \times \qty{e-9}{m} = \qty{7.50e-7}{m}$।
\emph{2.} $\qty{2.5}{km} = \qty{2.5e3}{m}$।
\emph{3.} $\qty{3.2}{\micro m} = \qty{3.2e-6}{m}$।
\emph{4.} $\qty{0.15}{Tm} = 0.15 \times \qty{e12}{m} = \qty{1.5e11}{m}$ --- यानी एक \omterm{def:g10:orders-of-magnitude:lightyear}{खगोलीय मात्रक}।
\end{solution}

\begin{solution}{exo:g10:orders-of-magnitude:2}
\emph{1.} $\qty{149600000}{km} = \qty{1.496e11}{m}$: $4$ \omterm{def:g10:orders-of-magnitude:sigfig}{सार्थक अंक}।
\emph{2.} $\qty{1.06e-10}{m}$: $3$ \omterm{def:g10:orders-of-magnitude:sigfig}{सार्थक अंक} (आरंभ के शून्य कभी नहीं
गिने जाते)।
\emph{3.} $\num{8.1e9}$: $2$ \omterm{def:g10:orders-of-magnitude:sigfig}{सार्थक अंक}।
\emph{4.} $\qty{5.20e-3}{s}$: $3$ \omterm{def:g10:orders-of-magnitude:sigfig}{सार्थक अंक} (अंत का शून्य एक वादा है)।
\end{solution}

\begin{solution}{exo:g10:orders-of-magnitude:3}
अपूर्णांशों और घातांकों को अलग-अलग लीजिए (\cref{prop:g10:orders-of-magnitude:powers}):
\[
3 \times 2 \times 10^{8-5} = \num{6e3}, \qquad
\tfrac{8}{4} \times 10^{4-(-3)} = \num{2e7},
\]
\[
5^{2} \times 10^{12} = \num{2.5e13}, \qquad
\tfrac{4 \times 3}{6} \times 10^{-3-9+7} = \num{2e-5} .
\]
\end{solution}

\begin{solution}{exo:g10:orders-of-magnitude:4}
मनुष्य: $1.7 < 5$, कोटि $\qty{e0}{m}$। एवरेस्ट: $\qty{8848}{m} = \qty{8.848e3}{m}$ और $8.848 \geq 5$: कोटि $\qty{e4}{m}$।
जीवाणु: कोटि $\qty{e-6}{m}$। पृथ्वी: $5.97 \geq 5$, कोटि $\qty{e25}{kg}$।
\end{solution}

\begin{solution}{exo:g10:orders-of-magnitude:5}
\emph{1.} आधा अंशांक: $L = (127.0 \pm 0.5)\,\unit{mm}$।
\emph{2.} $\Delta L / L = 0.5/127 \approx 0.4\%$।
\emph{3.} रबड़: $0.5/12 \approx 4\%$। पेंसिल वाली माप लगभग दस गुना अधिक परिशुद्ध है
--- वही मापनी, वही \omterm{def:g10:orders-of-magnitude:uncertainty}{निरपेक्ष अनिश्चितता}, पर वस्तु दस गुना लंबी।
\end{solution}

\begin{solution}{exo:g10:orders-of-magnitude:6}
\emph{1.} प्रॉक्सिमा: $4.24 \times \qty{9.47e15}{m} = \qty{4.01e16}{m}$। सिरियस: $8.6 \times \qty{9.47e15}{m} = \qty{8.1e16}{m}$ (दो \omterm{def:g10:orders-of-magnitude:sigfig}{सार्थक अंक},
जितने आगत में हैं)।
\emph{2.} $t = d/c = \qty{2.4e13}{m} / \qty{3.00e8}{m/s}
= \qty{8.0e4}{s}$, यानी $\num{8.0e4}/3600 \approx \qty{22}{h}$: यान के हर संदेश के लिए लगभग पूरा एक दिन।
\end{solution}

\begin{solution}{exo:g10:orders-of-magnitude:7}
हर बार $t = d/c$। चंद्रमा: $\qty{3.84e8}{m}/\qty{3.00e8}{m/s} = \qty{1.28}{s}$। सूर्य: $\qty{1.496e11}{m}/\qty{3.00e8}{m/s} = \qty{499}{s} \approx
\qty{8.3}{min}$। नेपच्यून: $\qty{4.50e12}{m}/\qty{3.00e8}{m/s} = \qty{1.50e4}{s} \approx
\qty{4.2}{h}$। जब आप
नेपच्यून को देखते हैं, तो आप उसे लगभग चार घंटे पुराना देखते हैं --- हर
दूरबीन एक हल्की-फुल्की समय-मशीन है।
\end{solution}

\begin{solution}{exo:g10:orders-of-magnitude:8}
\emph{1.} एक शीट: $\qty{52.0}{mm}/500 = \qty{0.104}{mm}
= \qty{1.04e-4}{m}$। अनिश्चितता भी उसी तरह बँटती है: $\qty{0.5}{mm}/500 = \qty{0.001}{mm}$,
इसलिए $t = (104 \pm 1)\,\unit{\micro m}$।
\emph{2.} $1/104 \approx 1\%$।
\emph{3.} एक अकेली शीट को सीधे मापने पर $(0.1 \pm 0.5)\,\unit{mm}$ मिलता: लगभग $500\%$ की
\omterm{def:g10:orders-of-magnitude:uncertainty}{सापेक्ष अनिश्चितता} --- मापनी तो यह भी सिद्ध नहीं कर पाती कि शीट
है भी। $500$ शीटों को ढेर करने से मापी गई लंबाई $500$ गुना हो जाती
है, जबकि मापनी की \omterm{def:g10:orders-of-magnitude:uncertainty}{निरपेक्ष अनिश्चितता} जस की तस रहती है: इसलिए
\omterm{def:g10:orders-of-magnitude:uncertainty}{सापेक्ष अनिश्चितता} उसी $500$ गुणक से सिकुड़ जाती है।
\end{solution}

\begin{solution}{exo:g10:orders-of-magnitude:9}
हर लंबाई को $\qty{2e-10}{m}$ से भाग दीजिए। जीवाणु: $\qty{2.0e-6}{m}/\qty{2e-10}{m} = \num{1.0e4}$ परमाणु। बाल: $\qty{7.0e-5}{m}/\qty{2e-10}{m} = \num{3.5e5}$।
मनुष्य: $1.7/\num{2e-10} = \num{8.5e9}$, \omterm{def:g10:orders-of-magnitude:oom}{परिमाण की कोटि} $10^{10}$: सिर से पैर तक दस अरब
परमाणु।
\end{solution}

\begin{solution}{exo:g10:orders-of-magnitude:10}
\emph{1.} $s = \qty{0.22}{m}/\qty{1.39e9}{m} = \num{1.58e-10}$।
\emph{2.} पृथ्वी: $\qty{1.27e7}{m} \times \num{1.58e-10} = \qty{2.0e-3}{m} =
\qty{2.0}{mm}$ --- यानी काली मिर्च का एक दाना --- जो फुटबॉल
से $\qty{1.496e11}{m} \times \num{1.58e-10} = \qty{23.7}{m} \approx
\qty{24}{m}$ दूर परिक्रमा करता है।
\emph{3.} चंद्रमा: उस दाने से $\qty{3.84e8}{m} \times \num{1.58e-10} = \qty{6.1}{cm}$ दूर। नेपच्यून: $\qty{4.50e12}{m} \times \num{1.58e-10} \approx \qty{710}{m}$: प्रतिरूप
स्कूल के मैदान में नहीं समाता।
\end{solution}

\begin{solution}{exo:g10:orders-of-magnitude:11}
\emph{1.} $21.0 \times 29.7 = \qty{623.7}{cm^2}$, जिसे आगतों के तीन सार्थक अंकों तक पूर्णांकित
करने पर $\qty{624}{cm^2}$।
\emph{2.} $v = \qty{100.0}{m}/\qty{12.3}{s} = \qty{8.13}{m/s}$ (तीन अंक, जो समय ने तय किए)।
\emph{3.} पाँच \omterm{def:g10:orders-of-magnitude:sigfig}{सार्थक अंक} क्षेत्रफल को $\qty{0.01}{cm^2}$ तक जानने का दावा
करते, जबकि मिलीमीटर स्तर के आगतों ने इतनी परिशुद्धि कभी दी ही नहीं ---
गढ़ा हुआ दशमलव फिर भी जालसाज़ी ही है।
\end{solution}

\begin{solution}{exo:g10:orders-of-magnitude:12}
\emph{1.} त्रिज्या $r = \qty{0.55}{m}$, क्षेत्रफल $A = \pi r^{2} = \pi \times 0.55^{2} \approx \qty{0.95}{m^2}$। आयतन $V = \qty{1.0}{mm^3} = \qty{1.0e-9}{m^3}$, और परत
(क्षेत्रफल) $\times$ (मोटाई) होती है, इसलिए
\[
e = \frac{V}{A} = \frac{\qty{1.0e-9}{m^3}}{\qty{0.95}{m^2}}
= \qty{1.1e-9}{m} .
\]
\emph{2.} अणु की कोटि $\qty{e-9}{m}$ है --- परमाणु के $\qty{e-10}{m}$ से एक डंडा ऊपर,
जो इस बात से मेल खाता है कि अणु परमाणुओं का एक छोटा गुच्छा होता है।
फ्रैंकलिन के एक चम्मच तेल ने आधा एकड़ तालाब शांत कर दिया और, उन्हें
पता चले बिना, नैनोमीटर नाप दिया।
\end{solution}

\begin{solution}{exo:g10:orders-of-magnitude:13}
मान्यताएँ: विश्राम में लगभग $70$ धड़कन प्रति मिनट, और $80$ वर्ष की
आयु। प्रति दिन: $70 \times 60 \times 24 \approx \num{1e5}$ धड़कनें। प्रति वर्ष: $\num{1e5} \times 365 \approx \num{4e7}$। पूरे जीवन में:
$\num{4e7} \times 80 \approx \num{3e9}$ --- \omterm{def:g10:orders-of-magnitude:oom}{परिमाण की कोटि} $10^{9}$, यानी कुछ अरब धड़कनें। चार
\omterm{def:g10:orders-of-magnitude:sigfig}{सार्थक अंक} यह दिखावा करते कि अस्सी साल तक हृदय-गति मिनट-दर-मिनट
पता है: आगत तो \omterm{def:g10:orders-of-magnitude:oom}{परिमाण की कोटि} भर के अनुमान हैं, और कोई निर्गत अपने आगतों
से ऊपर नहीं जा सकता (\cref{met:g10:orders-of-magnitude:sigfig})।
\end{solution}

\begin{solution}{exo:g10:orders-of-magnitude:14}
$n$ तहों के बाद मोटाई $2^{n} \times \qty{1.0e-4}{m}$ होती है।
\emph{1.} चाहिए $2^{n} \geq 8848/\num{1.0e-4} = \num{8.8e7}$। चूँकि $2^{26} \approx \num{6.7e7}$ कम पड़ता है और $2^{27} \approx \num{1.3e8}$ काफ़ी है:
$27$ तहें (ढेर $\approx \qty{13}{km}$)।
\emph{2.} चाहिए $2^{n} \geq \num{3.84e12}$: $2^{41} \approx
\num{2.2e12}$ कम है, $2^{42} \approx \num{4.4e12}$ पहुँच जाता है ---
$42$ तहें, ढेर $\qty{4.4e8}{m}$, यानी चंद्रमा से आराम से आगे।
\emph{3.} केवल प्रयोग गलत होता है: हर तह क्षेत्रफल आधा कर देती है, और
सात तहों के बाद ढेर अपनी चौड़ाई से अधिक मोटा हो जाता है। दुगुना करने का
अंकगणित अछूता रहता है --- चरघातांकी वृद्धि बस कागज़ (और अंतर्ज्ञान) की
पहुँच से आगे निकल जाती है।
\end{solution}

\begin{solution}{exo:g10:orders-of-magnitude:15}
\emph{1.} चंद्रमा: $\qty{0.03}{m}/\qty{3.84e8}{m} = \num{7.8e-11}$। मेज़: $\qty{0.001}{m}/\qty{1.000}{m} = \num{1.0e-3}$। चंद्रमा वाली माप $\num{1.0e-3}/\num{7.8e-11} \approx \num{1.3e7}$
गुना अधिक परिशुद्ध है --- यानी एक करोड़ से भी अधिक।
\emph{2.} एक \omterm{def:g10:orders-of-magnitude:unit}{मीटर} की मेज़ पर $\num{7.8e-11}$ की बराबरी करने का अर्थ है
$\Delta x = \qty{7.8e-11}{m}$, जो एक परमाणु ($\approx \qty{e-10}{m}$) से भी छोटा है: मेज़ का किनारा तो इतनी
पैनी तरह \emph{परिभाषित} भी नहीं है। सापेक्ष रूप से देखें तो
पृथ्वी--चंद्रमा की लेज़र दूरी हमारी प्रजाति की सबसे बारीक मापों में से
एक है।
\end{solution}

\begin{solution}{pb:g10:orders-of-magnitude:1}
\textbf{1.} कण: $\qty{0.5}{mm} = \qty{5e-4}{m}$, यानी मिली पैमाने पर; परमाणु: $\qty{2.5e-10}{m} = \qty{0.25}{nm}$, यानी
नैनो पैमाने पर।

\textbf{2.} $\qty{5e-4}{m}/\qty{2.5e-10}{m} = \num{2.0e6}$: आरपार बीस लाख परमाणु।

\textbf{3.} हर भुजा पर $2 \times 10^{6}$ परमाणुओं वाला घन $(2\times10^{6})^{3} = 8 \times 10^{18}$ परमाणु समेटता
है: \omterm{def:g10:orders-of-magnitude:oom}{परिमाण की कोटि} $10^{19}$।

\textbf{4.} $V = (\qty{5e-4}{m})^{3} = \qty{1.25e-10}{m^3}$, इसलिए $m = \rho V = \qty{2500}{kg/m^3} \times \qty{1.25e-10}{m^3}
= \qty{3.1e-7}{kg} = \qty{0.31}{mg}$।

\textbf{5.} प्रति परमाणु द्रव्यमान: $\qty{3.1e-7}{kg}/(8\times10^{18}) = \qty{3.9e-26}{kg}$ --- ठीक ``कुछ $10^{-26}\,\unit{kg}$''
वाली सीमा में। परमाणुओं के इस भोले घन से असली परमाणु-द्रव्यमान निकल आते
हैं: प्रतिरूप जाँच (\cref{met:g10:orders-of-magnitude:sanity}) में पूरे अंकों से पास होता है।

\textbf{6.} $s = \qty{5e-4}{m}/\qty{1.39e9}{m} = \num{3.6e-13}$।

\textbf{7.} प्रतिरूप-पृथ्वी: व्यास $\qty{1.27e7}{m} \times \num{3.6e-13} = \qty{4.6e-6}{m}$ --- यानी एक जीवाणु जितनी ---
और कण से $\qty{1.496e11}{m} \times \num{3.6e-13} = \qty{5.4}{cm}$ की दूरी पर: पृथ्वी रेत के एक कण से पाँच सेंटीमीटर दूर
एक सूक्ष्मजीव है।

\textbf{8.} $\qty{4.5e12}{m} \times \num{3.6e-13} = \qty{1.6}{m}$: पूरा ग्रह-मंडल एक खाने की मेज़ के भीतर परिक्रमा
करता है।

\textbf{9.} $\qty{4.0e16}{m} \times \num{3.6e-13} = \qty{1.4e4}{m}
= \qty{14}{km}$।

\textbf{10.} हमारे मेज़ भर के सौर मंडल का सबसे पास का तारा चौदह
किलोमीटर दूर रेत का एक और कण है, और बीच में लगभग कुछ भी नहीं: आकाशगंगा
भारी बहुमत से खाली जगह है --- यही कारण है कि तारों की टक्कर लगभग
अनसुनी बात है।

\textbf{11.} $t = \qty{5e-4}{m}/\qty{3.00e8}{m/s} = \qty{1.7e-12}{s}$ --- यानी दो पिकोसेकंड से भी कम।

\textbf{12.} चंद्रमा: $\qty{3.84e8}{m}/\qty{3.00e8}{m/s} =
\qty{1.28}{s}$। सूर्य: $\qty{1.496e11}{m}/\qty{3.00e8}{m/s} =
\qty{499}{s} \approx \qty{8.3}{min}$: धूप हमेशा आठ मिनट पुरानी
होती है।

\textbf{13.} एक वर्ष $365.25 \times 86400 = \qty{3.156e7}{s}$ का होता है, इसलिए $\qty{1}{ly} = \qty{3.00e8}{m/s} \times \qty{3.156e7}{s}
= \qty{9.47e15}{m}$; फिर $\qty{4.0e16}{m}/\qty{9.47e15}{m} = 4.2$:
आँकड़ों वाली दूरी सचमुच लगभग $\qty{4.2}{ly}$ है।

\textbf{14.} $\qty{2.4e22}{m}/\qty{9.47e15}{m} = \num{2.5e6}$: पच्चीस लाख \omterm{def:g10:orders-of-magnitude:lightyear}{प्रकाश वर्ष}। वह प्रकाश
एंड्रोमेडा से लगभग $2.5$ लाख वर्ष पहले चला था, जब \emph{होमो} वंश के
पहले सदस्य पत्थर गढ़ना सीख रहे थे --- और \emph{होमो सेपियन्स} तब कहीं
था ही नहीं।

\textbf{15.} $\qty{3.00e8}{m/s} \times \qty{1e-9}{s} = \qty{0.30}{m}$। एक-तिहाई नैनोसेकंड में कोई संकेत अधिक से अधिक
लगभग $\qty{10}{cm}$ तय करता है, इसलिए मशीन के जिन दो हिस्सों को एक ही घड़ी-टिक
के भीतर बात करनी है उन्हें एक-दूसरे से दस सेंटीमीटर के भीतर बैठना पड़ता
है: चाल आकार पर सीमा लगा देती है।

\textbf{16.} पानी में औसत परमाणु: $\qty{3.0e-26}{kg}/3 = \qty{1.0e-26}{kg}$। परमाणुओं की संख्या: $\qty{70}{kg}/\qty{1.0e-26}{kg} = \num{7e27}$।

\textbf{17.} कण के सामने: $\num{7e27}/\num{8e18} \approx 10^{9}$ --- आप में एक अरब कणों जितने परमाणु
हैं। तारों के सामने: $\num{7e27}/\num{2e23} = \num{3.5e4}$ --- आपके परमाणु प्रेक्षणीय ब्रह्मांड के
सारे तारों से पैंतीस हज़ार गुना अधिक हैं।

\textbf{18.} $\text{(आप)}/\text{(परमाणु)} =
\qty{1.7}{m}/\qty{2.5e-10}{m} \approx \num{7e9} \approx 10^{10}$, जबकि $\text{(ब्रह्मांड)}/\text{(आप)} =
\qty{8.8e26}{m}/\qty{1.7}{m} \approx \num{5e26}$। सीढ़ी पर आप परमाणु से लगभग दस डंडे
ऊपर खड़े हैं और ब्रह्मांड से लगभग सत्ताईस नीचे: गुणात्मक रूप से मनुष्य
अपने परमाणुओं के बहुत पास रहते हैं।

\textbf{19.} शर्त से $x^{2} = 10^{-15} \times 10^{26.9} = 10^{11.9}$ मिलता है, इसलिए $x = 10^{5.95} \approx \qty{9e5}{m} \approx \qty{900}{km}$ --- जो लगभग ठीक
सबसे बड़े क्षुद्रग्रह सीरीस के व्यास ($\approx \qty{940}{km}$) जितना है। नाभिक और
ब्रह्मांड के बीच सीढ़ी के ठीक आधे पर एक मामूली बौना ग्रह बैठा है।

\textbf{20.} $\qty{8.8e26}{m} \times \num{3.6e-13} = \qty{3.2e14}{m}$, यानी $\qty{3.2e14}{m}/\qty{1.496e11}{m} \approx \num{2.1e3}$ \omterm{def:g10:orders-of-magnitude:lightyear}{खगोलीय मात्रक}। और अंतिम चोट:
ब्रह्मांड को तब तक सिकोड़िए जब तक सूर्य रेत का कण न बन जाए, फिर भी
प्रेक्षणीय ब्रह्मांड असली पृथ्वी--सूर्य दूरी से दो हज़ार गुना फैला रहता
है --- कोई मापनी-प्रतिरूप ब्रह्मांड को मानवीय आकार तक नहीं ला पाता;
उसे केवल दस की घातें ढो सकती हैं।
\end{solution}
